% Track V candidate section. Compile with the main report preamble; no new packages.
\section{Commanded vertical motion}
\label{sec:v-vertical}

\paragraph{Scope.}
A quadrotor's altitude is part of every action-conditioned aerial world model's control
signal, and it is the one flatness axis that changes what the camera sees without changing the
thrust direction: a commanded climb from $z_0$ to $z_0+\Delta z$ must change the camera's height
above the ground by the ratio $(z_0+\Delta z)/z_0$. EyeHeight~\cite{eyeheight2026} measures
whether a world model \emph{keeps} the camera height of its start frame on constant-height
indoor trajectories; this section asks the complementary question, whether a commanded climb or
descent is \emph{depicted} with the commanded ratio, and whether the measurement can be
certified on real drone flights.

\paragraph{Estimator and statistic.}
Following EyeHeight we reconstruct every fourth frame in a single multi-view pass (Pi3X, the
successor of their Pi3, with the camera intrinsics supplied) in overlapping 8\,s chunks at
7.5\,Hz, fit a ground plane to the reconstruction's own point cloud, and read the camera's
signed distance $h(t)$ to that plane. The check statistic is
\begin{equation}
e = \log\frac{h(t_2)}{h(t_1)} - \log\frac{z(t_2)}{z(t_1)}
\label{eq:v-height-ratio}
\end{equation}
over a 3\,s window on the simulator (the scene ground is $z=0$ and the truth is exact) and
\begin{equation}
e = \log\frac{h(t_2)}{h(t_1)} - \log\left(1 + \frac{z(t_2)-z(t_1)}{s\,h(t_1)}\right)
\label{eq:v-height-ratio-real}
\end{equation}
on real flights, where $s$ is one similarity scale per reconstruction chunk (the ratio of RMS
radii of the camera-centre cloud to the Leica track) and $z$ is the truth altitude. The height
ratio cancels the reconstruction's scale; it does not cancel an error in the plane's
\emph{orientation}, which enters as $\delta \cdot L$ for a horizontal displacement $L$ and a
normal error $\delta$. A scale-free alternative, $\rho = \Delta z / L$, is reported with it.

\paragraph{Where the population is.}
Of the 200 development clips of the simulated set, 28 have commanded vertical motion: all of
them \texttt{climb\_turn}, climbing 9.4--48.4\,m from a 12\,m start (a per-window ratio of
1.13--1.31). The real pilot is UZH-FPV's three outdoor forward sequences; TII is excluded
because its fisheye distortion model is not carried by the reader feeds.

\paragraph{Result.}
On the rendered clips the estimator follows the commanded ratio: the median $|e|$ is
0.094 across 761 clips with 97.3\% reconstructable windows, and the
depicted height tracks the truth altitude per chunk. On the real outdoor flights it does not:
the pooled median $|e|$ is 0.480 over 95 windows from three sequences
(0.618, 0.502 and 0.433), and the per-chunk correlation between depicted height and truth altitude
ranges from -0.93 to 1.00. The implied ground-normal error is
7.1$^\circ$ (median over chunks, 1.3--53.8), against the
$\sim0.5^\circ$ that a $h \approx 2$\,m camera over a $L \approx 20$\,m baseline would need
(per-sequence medians 7.1$^\circ$, 11.2$^\circ$ and 24.9$^\circ$); the same quantity on the
rendered clips is 6.6$^\circ$, so the plane is fitted about equally badly in both domains and what
differs is the leverage $L/h$ (0.23--1.63 rendered, 7.9--12.0 real).
The pre-registered stop rule compared the real error with the effect of the smallest labelled
vertical defect (a $\times 0.5$ or $\times 1.5$ scaling of a per-window ratio near 1.2, i.e.\
$|\log|$ of 0.06--0.11): the real error exceeds it on 28 of 28 vertical-motion clips, so
the check stops at feasibility and no defect render, real false-reject table or generated-clip
arm is reported. The vertex of the failure is geometric, not a threshold: at
$h \sim 1$--3\,m over $L \sim 15$--45\,m of travel, the plane's orientation must be known to a
few tenths of a degree, and a bumpy real ground does not provide it.

\paragraph{What this does and does not show.}
It shows that EyeHeight's measurement transfers to \emph{commanded} vertical motion on rendered
clips at 12--60\,m altitude, and that it does not transfer to low, fast drone footage anywhere among the
14 UZH rigs tried. It does not show that no vertical-motion check is possible: a
near-field-only plane fit, a gravity-referenced height from inertial-aligned visual odometry, or
a plane-free comparison of depicted vertical velocity against a metrically scaled path would all
change the estimator, and none was tested. A downward-looking rig was tested and does not help
(the 45$^\circ$-down UZH sequences have median $|e|$ 0.37--0.45 because their flights are low,
$L/h \approx 6$--7). The
generated-clip arm is unrun, so nothing is claimed about whether a particular generator follows
a commanded climb.

\begin{table}[t]
\centering
\small
\begin{tabular}{@{}lrrrr@{}}
\toprule
 & windows & median $|e|$ & p90 $|e|$ & covered \\
\midrule
Simulated, 28 climb clips (rendered) & 761/782 & 0.094 & 0.227 & 97.3 \\
UZH outdoor forward, 1 & 45/54 & 0.618 & 1.453 & 0.83 \\
UZH outdoor forward, 3 & 33/54 & 0.502 & 1.816 & 0.61 \\
UZH outdoor forward, 5 & 17/24 & 0.433 & 1.600 & 0.71 \\
\bottomrule
\end{tabular}
\caption{Commanded-vertical-motion feasibility. $e$ is Eq.~\ref{eq:v-height-ratio} on the
simulator and Eq.~\ref{eq:v-height-ratio-real} on the real flights; windows are 3\,s with a
0.5\,s stride. Coverage is the fraction of windows whose window is reconstructable.}
\label{tab:v-feasibility}
\end{table}
